% =========================================================
% Stage 1 S1-A
% =========================================================
\part{Stage 1 --- 선형 데이터 처리와 Associative Containers}
\chapter{S1-A --- Associative Containers}

% =========================================================
% Part A
% =========================================================
\section{Part A --- 위치와 선수지식}

\subsection{현재 위치}
Stage 0에서 기본 C++ 실행, container, iterator, reference, 복잡도와 numeric safety를 다뤘다면, S1-A에서는 값을 단순히 순서대로 저장하는 것을 넘어 \textbf{key를 이용해 찾고, 순서를 유지하고, 중복을 보존하거나 제거하는 자료구조}를 다룬다.

이번 Learning Unit의 핵심 범위는 다음과 같다.
\begin{itemize}
\item \code{map}, \code{unordered\_map}
\item \code{set}, \code{unordered\_set}
\item \code{multiset}
\item frequency counting과 key $\to$ value 상태 관리
\item \code{find()}, \code{erase()}, \code{operator[]}
\item ordered container의 \code{lower\_bound()}, \code{upper\_bound()}
\item predecessor / successor와 iterator boundary
\item ordered vs unordered complexity
\item 문자열 key에서의 aggregate payload $L$
\item forward index와 reverse index
\end{itemize}

\subsection{왜 지금 배우는가}
Competitive Programming에서는 다음 요구가 매우 자주 등장한다.
\begin{itemize}
\item 값이 이전에 등장했는지 빠르게 확인한다.
\item 어떤 key에 대응하는 count, score, balance 등을 갱신한다.
\item 현재 저장된 값 중 최소/최대 또는 $x$ 주변 값을 찾는다.
\item 중복 원소를 그대로 보존하면서 정렬 상태를 유지한다.
\item 동적으로 삽입/삭제되는 데이터에 query를 반복한다.
\end{itemize}

이때 모든 문제를 \code{vector}와 선형 탐색으로 해결하려 하면 $O(N^2)$으로 커지는 경우가 많다. Associative container의 핵심은 API를 외우는 것이 아니라 \textbf{요구 연산을 보고 적절한 container를 선택하는 것}이다.

\subsection{Core Decision Boundary}
S1-A의 핵심 선택 경계는 다음과 같다.
\begin{center}
\begin{tabularx}{\textwidth}{p{4.2cm}X}
\toprule
선택 상황 & 판단 기준 \\
\midrule
key만 필요한가, key $\to$ value가 필요한가? & membership만 필요하면 set 계열, key마다 별도 상태가 필요하면 map 계열을 우선 고려한다. \\
ordering이 필요한가? & 정렬 순서, 최소/최대, predecessor/successor, 범위 경계가 필요하면 ordered container를 고려한다. 순서가 필요 없고 평균적인 빠른 lookup이 중요하면 unordered 계열을 고려한다. \\
중복을 보존해야 하는가? & 동일 값 여러 개를 독립적으로 유지해야 하면 \code{multiset} 같은 duplicate-preserving 구조가 필요하다. \\
동적 update인가, 정적 데이터인가? & 삽입/삭제가 반복되는가, 아니면 한 번 정렬한 뒤 읽기만 하는가를 구분한다. 정적 데이터라면 정렬된 \code{vector}가 더 단순할 수 있다. \\
한 방향 조회만 필요한가? & $A\to B$와 $B\to A$ query가 모두 자주 등장하면 reverse index를 별도로 유지할지 검토한다. \\
\bottomrule
\end{tabularx}
\end{center}

\subsection{Immediate Coverage Floor}
이 Unit의 즉시 학습 증거는 다음을 포함해야 한다.
\begin{itemize}
\item key $\to$ value를 다루는 map 계열 구현
\item 중복을 보존하는 ordered container 구현
\item ordered / unordered 중 하나를 고르는 이유를 complexity와 요구 순서로 설명
\item 존재하지 않는 key에 대한 \code{find()}, \code{operator[]}의 차이를 설명
\end{itemize}

% =========================================================
% Part B
% =========================================================
\section{Part B --- 학습 목표}
이 Unit을 마치면 다음 행동을 독립적으로 수행할 수 있어야 한다.
\begin{enumerate}
\item membership 문제와 key $\to$ value 문제를 구분한다.
\item \code{map}과 \code{unordered\_map}의 ordering과 complexity 차이를 설명한다.
\item \code{set}, \code{unordered\_set}, \code{multiset}의 중복 보존 차이를 설명한다.
\item \code{operator[]}가 존재하지 않는 key를 삽입할 수 있다는 점을 이해한다.
\item \code{find()} 결과를 \code{end()}와 비교한 뒤 안전하게 dereference한다.
\item \code{erase(key)}와 \code{erase(iterator)}의 차이를 상황에 맞게 사용한다.
\item ordered container에서 최소/최대와 이웃 원소를 효율적으로 찾는다.
\item \code{lower\_bound()}, \code{upper\_bound()}의 의미를 구분한다.
\item \code{prev(begin())}, \code{*end()}가 왜 잘못된지 설명한다.
\item hash container의 expected complexity와 worst-case complexity를 구분한다.
\item 문자열 key 문제에서 원소 수뿐 아니라 총 문자 수 $L$이 complexity에 영향을 줄 수 있음을 인식한다.
\item 큰 mapped container를 실수로 복사하지 않도록 \code{const auto\&}를 사용할 수 있다.
\item query 방향이 반대인 경우 reverse index가 필요한지 판단한다.
\end{enumerate}

% =========================================================
% Part C
% =========================================================
\section{Part C --- 개념과 원리}

\subsection{C1. Associative container의 기본 생각}
Sequence container는 보통 위치를 기준으로 원소를 다룬다.
\begin{lstlisting}
vector<int> v = {10, 20, 30};
cout << v[1];
\end{lstlisting}

Associative container는 key 또는 값 자체를 기준으로 검색한다.
\begin{lstlisting}
map<string, int> price;
price["BTC"] = 100;

set<int> active;
active.insert(7);
\end{lstlisting}

따라서 문제를 읽을 때 ``몇 번째 원소인가?''보다 ``어떤 key인가?'', ``이 값이 존재하는가?'', ``이 값의 바로 왼쪽/오른쪽은 무엇인가?''가 중요하다면 associative container를 후보에 올린다.

\subsection{C2. \code{map}: ordered key $\to$ value}
\code{map<K,V>}는 key마다 value 하나를 대응시키며 key 순서를 정렬된 상태로 유지한다.
\begin{lstlisting}
map<string, int> balance;
balance["BTC"] = 120;
balance["ETH"] = 70;

for (const auto& [symbol, value] : balance) {
    cout << symbol << ' ' << value << '\n';
}
\end{lstlisting}

대표 연산은 보통 $O(\log N)$이다.
\begin{itemize}
\item insert / assignment by key: $O(\log N)$
\item \code{find}: $O(\log N)$
\item \code{erase(key)}: $O(\log N)$에 탐색 비용이 포함된다.
\item \code{begin()}: 정렬상 첫 key에 접근한다.
\end{itemize}

\subsection{C3. \code{unordered\_map}: ordering 없이 평균적으로 빠른 key $\to$ value}
\code{unordered\_map<K,V>}는 hash table을 사용한다. 정렬 순서를 보장하지 않는 대신 lookup, insert, erase가 평균적으로 $O(1)$이다.
\begin{lstlisting}
unordered_map<string, int> freq;
++freq["apple"];
++freq["apple"];
++freq["pear"];
\end{lstlisting}

하지만 worst case에서는 많은 key가 같은 bucket에 몰릴 수 있다. $K$개의 key가 저장되어 있을 때 하나의 lookup이 $O(K)$까지 나빠질 수 있고, $Q$개의 연산 전체는 이론적으로 $O(Q^2)$까지 갈 수 있다.

\begin{keybox}[expected와 worst-case를 분리해서 쓰기]
\code{unordered\_map}, \code{unordered\_set}을 사용했다면 ``$O(1)$''이라고만 쓰지 말고, 정식 complexity 분석에서는 필요한 경우 \textbf{expected}와 \textbf{worst-case}를 구분한다.
\end{keybox}

\subsection{C4. \code{set}: ordered unique values}
\code{set<T>}는 중복 없이 값을 정렬된 상태로 저장한다.
\begin{lstlisting}
set<int> cuts = {0, 100};
cuts.insert(30);
cuts.insert(70);
\end{lstlisting}

같은 값을 다시 삽입해도 원소 수는 증가하지 않는다.
\begin{lstlisting}
cuts.insert(30); // 30은 하나만 유지
\end{lstlisting}

membership뿐 아니라 최소/최대, predecessor/successor query가 필요할 때 강력하다.

\subsection{C5. \code{unordered\_set}: ordering 없는 unique membership}
값의 존재 여부만 필요하고 정렬 순서는 필요 없다면 \code{unordered\_set}이 자연스럽다.
\begin{lstlisting}
unordered_set<int> seen;
if (seen.find(x) == seen.end()) {
    seen.insert(x);
}
\end{lstlisting}

평균적으로 lookup과 insert는 $O(1)$이지만 hash collision에 의한 worst case는 별도로 존재한다.

\subsection{C6. \code{multiset}: ordered duplicate-preserving values}
\code{multiset<T>}는 \code{set}과 달리 같은 값을 여러 번 저장할 수 있다.
\begin{lstlisting}
multiset<int> values;
values.insert(5);
values.insert(5);
values.insert(2);
\end{lstlisting}

이때 저장 상태는 개념적으로 \code{2, 5, 5}이다.

한 occurrence만 삭제하고 싶다면 iterator를 먼저 찾은 뒤 그 iterator를 삭제한다.
\begin{lstlisting}
auto it = values.find(5);
if (it != values.end()) {
    values.erase(it);
}
\end{lstlisting}

\subsection{C7. \code{operator[]}는 단순 조회가 아니다}
map 계열에서 다음 코드는 매우 편리하다.
\begin{lstlisting}
++freq[key];
\end{lstlisting}

하지만 \code{operator[]}는 key가 없으면 default value를 가진 새 원소를 삽입한다.
\begin{lstlisting}
map<string, int> mp;
cout << mp["missing"]; // 0 출력, 동시에 key가 삽입됨
\end{lstlisting}

따라서 ``존재 여부만 확인''하는 query에서 \code{mp[x]}를 쓰면 container의 상태와 공간 사용량이 바뀔 수 있다.

\subsection{C8. \code{find()}와 \code{end()}: 안전한 lookup}
\begin{lstlisting}
auto it = mp.find(key);
if (it != mp.end()) {
    cout << it->second << '\n';
}
\end{lstlisting}

\code{find()}는 key가 없으면 \code{end()}를 반환한다. \code{end()}는 실제 원소를 가리키지 않으므로 먼저 비교한 뒤에 dereference해야 한다.

C++20에서는 associative container에 \code{contains()}도 사용할 수 있다.
\begin{lstlisting}
if (s.contains(x)) {
    // ...
}
\end{lstlisting}

단, 정확히 C++17만 허용되는 환경에서는 \code{contains()}를 사용할 수 없다. ``C++17 이상''처럼 C++20도 허용된다면 컴파일 표준을 C++20으로 선택할 수 있다.

\subsection{C9. \code{erase}: key를 지울지, iterator를 지울지}
container와 overload에 따라 의미와 비용이 다르다.
\begin{lstlisting}
multiset<int> ms = {2, 5, 5, 5};

auto it = ms.find(5);
if (it != ms.end()) {
    ms.erase(it); // 한 occurrence
}
\end{lstlisting}

반면 \code{ms.erase(5)}는 같은 key의 모든 occurrence를 지운다. 문제에서 ``하나만 삭제''인지 ``전부 삭제''인지 반드시 확인해야 한다.

\subsection{C10. ordered와 unordered의 iteration 순서}
\code{map}, \code{set}, \code{multiset}은 정렬 순서대로 iteration된다. 따라서 lexicographically smallest key가 필요하다면 \code{map.begin()}을 사용할 수 있다.

반면 \code{unordered\_map}, \code{unordered\_set}의 iteration 순서는 문제 해결에 사용할 수 있는 정렬 순서가 아니다.

\subsection{C11. 문자열 key에서는 원소 수만 세면 부족할 수 있다}
문자열을 hash하거나 비교하려면 문자열의 문자를 읽어야 한다. $N$개의 update 문자열과 $M$개의 query 문자열의 총 문자 수를 $L$이라 두면, hash 기반 알고리즘의 expected time을 단순히 $O(N+M)$이라고만 적기보다
\[
O(N+M+L)
\]
처럼 payload를 드러내는 편이 더 정확할 수 있다.

문자열 길이가 30 이하처럼 상수로 제한되어 있다면 $L=\Theta(N+M)$이므로 $O(N+M)$과 asymptotically 동치다. 그래도 $L$을 명시하면 어떤 비용이 숨겨져 있는지 더 분명하게 설명할 수 있다.

\subsection{C12. 정적 데이터와 동적 update를 구분하기}
정렬된 \code{vector}도 binary search를 이용해 빠르게 lookup할 수 있다. 그러나 중간 삽입/삭제가 반복되면 원소 이동 비용이 커질 수 있다.

\begin{center}
\begin{tabularx}{\textwidth}{p{3.7cm}XX}
\toprule
구조 & 강점 & 주의점 \\
\midrule
정렬된 \code{vector} & 메모리 locality, 정적 데이터에서 단순하고 빠름 & 중간 삽입/삭제가 $O(N)$ \\
\code{set}/\code{map} & 동적 insert/erase와 ordered lookup이 $O(\log N)$ & 상수 계수와 메모리 overhead가 큼 \\
\bottomrule
\end{tabularx}
\end{center}

\subsection{C13. Forward index와 reverse index}
데이터를
\[
A \to B
\]
방향으로만 저장했는데 query가 자주
\[
B \to A
\]
를 요구한다면 매번 전체 $A$를 훑게 될 수 있다.

예를 들어 box마다 card 목록을 저장한 경우,
\begin{lstlisting}
map<int, multiset<int>> box_to_cards;
\end{lstlisting}
만으로 ``card $x$가 들어 있는 모든 box''를 반복해서 찾으면 모든 box를 검사해야 한다.

이때 reverse index를 함께 유지할 수 있다.
\begin{lstlisting}
map<int, set<int>> card_to_boxes;
\end{lstlisting}

삽입 시 두 구조를 동시에 갱신한다.
\begin{lstlisting}
box_to_cards[box].insert(card);
card_to_boxes[card].insert(box);
\end{lstlisting}

update 비용과 공간은 늘지만 reverse query는 직접 접근할 수 있다.

\subsection{C14. \code{map::value\_type}과 숨은 container 복사}
\begin{lstlisting}
map<int, multiset<int>> boxes;
\end{lstlisting}
의 실제 element type은
\begin{lstlisting}
pair<const int, multiset<int>>
\end{lstlisting}
이다. key가 \code{const}인 이유는 tree의 ordering invariant를 임의로 깨뜨리지 못하게 하기 위해서다.

따라서 다음처럼 정확히 맞지 않는 pair type을 직접 적으면 변환 임시 객체가 생길 수 있다.
\begin{lstlisting}
for (const pair<int, multiset<int>>& p : boxes) {
    // ...
}
\end{lstlisting}

큰 mapped object의 숨은 복사를 피하려면 보통 다음이 안전하다.
\begin{lstlisting}
for (const auto& p : boxes) {
    // ...
}
\end{lstlisting}

또는 C++17 structured binding을 사용할 수 있다.
\begin{lstlisting}
for (const auto& [box, cards] : boxes) {
    // ...
}
\end{lstlisting}

\subsection{C15. \code{lower\_bound()}와 \code{upper\_bound()}}
ordered associative container의 member function을 사용하면 tree 구조를 이용해 $O(\log N)$에 경계를 찾는다.

\begin{itemize}
\item \code{s.lower\_bound(x)}: $x$ 이상인 첫 원소
\item \code{s.upper\_bound(x)}: $x$보다 큰 첫 원소
\end{itemize}

예를 들어
\begin{lstlisting}
set<int> s = {3, 7, 12, 20};
\end{lstlisting}
에서
\begin{lstlisting}
s.lower_bound(7);  // 7
s.upper_bound(7);  // 12
s.lower_bound(9);  // 12
s.upper_bound(20); // end()
\end{lstlisting}

\begin{notebox}[S1-B와의 경계]
정렬된 \code{vector}에 대한 generic \code{lower\_bound}, binary search family 자체는 다음 Unit에서 더 체계적으로 다룬다. 여기서는 ordered associative container의 핵심 query operation으로서 member \code{lower\_bound}/\code{upper\_bound}를 익힌다.
\end{notebox}

\subsection{C16. predecessor와 successor}
정렬된 set에서 $x$ 주변 값을 찾고 싶다면 경계 iterator를 이용한다.
\begin{lstlisting}
auto right = s.lower_bound(x);
if (right != s.end()) {
    // *right = x 이상인 최소값
}
if (right != s.begin()) {
    auto left = prev(right);
    // *left = x보다 작은 최대값
}
\end{lstlisting}

$x$보다 큰 strict successor가 필요하면 \code{upper\_bound(x)}가 직접적이다.

\subsection{C17. iterator boundary: \code{end()}는 있지만 before-begin은 일반적으로 없다}
STL range는 일반적으로 반열린 구간
\[
[\text{begin},\text{end})
\]
으로 표현한다. \code{end()}는 마지막 원소 다음의 종료 경계를 나타내지만 실제 원소는 아니다.

따라서 nonempty ordered container에서
\begin{lstlisting}
auto it = s.end();
--it; // 마지막 원소: valid
\end{lstlisting}
은 가능하지만
\begin{lstlisting}
auto it = s.begin();
--it; // invalid
\end{lstlisting}
은 허용되지 않는다.

역방향 순회가 필요하면 \code{rbegin()}, \code{rend()}라는 별도의 reverse range를 사용한다.

% =========================================================
% Part D
% =========================================================
\section{Part D --- Worked Example}

\subsection{D1. 동적 단어 빈도와 사전순 첫 단어}
다음 연산을 처리한다고 생각하자.
\begin{itemize}
\item \code{ADD s}: 문자열 \code{s}의 빈도를 1 증가시킨다.
\item \code{COUNT s}: 현재 빈도를 출력한다.
\item \code{FIRST}: 현재 저장된 문자열 중 사전순으로 가장 빠른 문자열과 빈도를 출력한다.
\end{itemize}

\code{COUNT}만 필요하다면 \code{unordered\_map}이 후보지만 \code{FIRST}가 있으므로 key ordering이 실제 요구사항이다. 따라서 \code{map<string,int>}가 자연스럽다.

\begin{lstlisting}
#include <iostream>
#include <map>
#include <string>
using namespace std;

int main() {
    int Q;
    cin >> Q;

    map<string, int> freq;

    while (Q--) {
        string op;
        cin >> op;

        if (op == "ADD") {
            string s;
            cin >> s;
            ++freq[s];
        }
        else if (op == "COUNT") {
            string s;
            cin >> s;
            auto it = freq.find(s);
            cout << (it == freq.end() ? 0 : it->second) << '\n';
        }
        else if (op == "FIRST") {
            if (freq.empty()) {
                cout << "EMPTY\n";
            }
            else {
                cout << freq.begin()->first << ' '
                     << freq.begin()->second << '\n';
            }
        }
    }
}
\end{lstlisting}

\subsection{D2. Complexity와 선택 이유}
각 ordered map operation은 저장된 distinct key 수를 $K$라 할 때 대체로 $O(\log K)$이다. 따라서 전체 query 수가 $Q$라면 기본적으로 $O(Q\log Q)$로 둘 수 있다. 문자열 비교 비용을 더 일반적으로 추적하고 싶다면 key length까지 고려할 수 있다.

이 예제의 핵심은 \textbf{평균 lookup 속도만 보고 unordered container를 선택하지 않고, 실제 query가 ordering을 요구하는지 먼저 확인하는 것}이다.

% =========================================================
% Part E
% =========================================================
\section{Part E --- 중간평가 문제}

\subsection{E-1. Ticker Directory}
\begin{problembox}{중간평가 문제 E-1 --- Ticker Directory}
처음에는 비어 있는 ticker directory가 있다. 다음 종류의 명령을 처리한다.
\begin{itemize}
\item \code{SET s x}: ticker \code{s}의 값을 \code{x}로 설정한다. 이미 존재하면 갱신한다.
\item \code{GET s}: \code{s}가 존재하면 값을 출력하고, 없으면 \code{NONE}을 출력한다.
\item \code{ERASE s}: \code{s}가 존재하면 삭제한다. 없으면 아무 일도 하지 않는다.
\item \code{FIRST}: directory가 비어 있지 않다면 사전순으로 가장 빠른 ticker와 그 값을 출력하고, 비어 있으면 \code{NONE}을 출력한다.
\end{itemize}

\textbf{구현 및 제출 조건}
\begin{itemize}
\item C++17 이상을 사용한다.
\item 완성 코드, 자료구조 선택 이유, 시간복잡도, 공간복잡도, edge case, 직접 측정한 풀이 시간을 제출한다.
\end{itemize}

\textbf{제약 조건}
\begin{TextBlock}
1 <= N <= 200000
1 <= ticker 길이 <= 20
ticker는 영문 대문자로 구성된다.
-1000000000 <= x <= 1000000000
\end{TextBlock}
\end{problembox}

\begin{notebox}[세션 복원 상태]
세션 검색으로 문제의 핵심 요구, 자료구조 선택, 복잡도, T\_solve와 평가 기록은 복원되었다. 다만 현재 회수 가능한 장기 기록에는 당시 제출 코드 전문과 sample 입력 전체가 verbatim으로 남아 있지 않아, 아래 코드는 기록으로 확인되는 실제 제출 구조만 보존한다. 확인 가능한 sample 출력에는 \code{100}, \code{BTC 100}, \code{120}, \code{NONE}, \code{ETH 70}이 순서대로 포함되어 있었다.
\end{notebox}

\begin{answerbox}[학습자 제출 구조]
학습자는 \code{map<string,int>}를 사용했다.
\begin{lstlisting}
map<string, int> ticker;

// SET
ticker[s] = x;

// GET
auto it = ticker.find(s);
if (it == ticker.end()) {
    cout << "NONE\n";
}
else {
    cout << it->second << '\n';
}

// ERASE
ticker.erase(s);

// FIRST
if (ticker.empty()) {
    cout << "NONE\n";
}
else {
    cout << ticker.begin()->first << ' '
         << ticker.begin()->second << '\n';
}
\end{lstlisting}

\textbf{선택 이유}: key $\to$ value 관계가 필요하고 \code{FIRST}에서 사전순 최소 key가 필요하므로 ordered \code{map}을 사용한다.

\textbf{시간복잡도}: $O(N\log N)$.

\textbf{공간복잡도}: 최대 $N$개의 ticker를 저장하므로 $O(N)$.

\textbf{edge case}: 명령이 하나뿐이고 존재하지 않는 ticker를 \code{GET}하는 경우.

\textbf{T\_solve}: 19분 51초.
\end{answerbox}

% =========================================================
% Q&A after Part E
% =========================================================
\section{중간평가 이후 학습 질의응답}

\subsection{Q1. 왜 \code{unordered\_map}이 아니라 \code{map}인가?}
\code{GET}, \code{SET}만 있었다면 평균 $O(1)$ lookup의 \code{unordered\_map}도 강한 후보다. 하지만 \code{FIRST}가 lexicographically smallest key를 요구한다. \code{map}은 key가 정렬되어 있어 \code{begin()}으로 바로 최소 key에 접근할 수 있다.

\subsection{Q2. \code{mp[q]}와 \code{find(q)}는 같은 조회인가?}
아니다. \code{mp[q]}는 key가 없으면 default value를 가진 새 원소를 삽입한다. 반면 \code{find(q)}는 container 상태를 바꾸지 않고 존재 여부를 확인한다. query가 상태를 변경하면 안 되는 경우에는 이 차이가 중요하다.

\subsection{Q3. 문자열 key에서 왜 총 길이 $L$을 따로 생각하는가?}
문자열을 hash하거나 비교하려면 문자를 읽어야 한다. 문자열 개수만 $N$으로 세면 payload 비용이 숨을 수 있다. 총 문자열 길이를 $L$이라 두면 expected complexity를 $O(N+M+L)$처럼 더 구조적으로 설명할 수 있다. 문자열 길이가 상수 상한을 가지면 다시 $O(N+M)$으로 단순화할 수 있다.

% =========================================================
% Part F
% =========================================================
\section{Part F --- 최종 성취도 평가 문제}

\subsection{F-1. Live Value Pool}
\begin{problembox}{최종평가 문제 F-1 --- Live Value Pool}
처음에는 비어 있는 정수 multiset이 있다. 다음 query를 처리한다.
\begin{itemize}
\item \code{ADD x}: 값 \code{x}를 하나 추가한다.
\item \code{CHECK x}: 현재 \code{x}가 하나 이상 존재하면 \code{YES}, 아니면 \code{NO}를 출력한다.
\item \code{REMOVE x}: \code{x}가 존재하면 \textbf{한 occurrence만} 삭제한다. 없으면 아무 일도 하지 않는다.
\item \code{RANGE}: 비어 있지 않으면 현재 최솟값과 최댓값을 출력하고, 비어 있으면 \code{EMPTY}를 출력한다.
\end{itemize}

\textbf{구현 및 제출 조건}
\begin{itemize}
\item C++17 이상을 사용한다.
\item 중복 값은 독립적으로 보존한다.
\item 완전한 코드, 자료구조 선택 이유, 시간복잡도, 공간복잡도, edge case, 직접 측정한 풀이 시간을 제출한다.
\end{itemize}

\textbf{제약 조건}
\begin{TextBlock}
1 <= Q <= 200000
-1000000000 <= x <= 1000000000
\end{TextBlock}
\end{problembox}

\begin{notebox}[세션 복원 상태]
세션 검색으로 문제의 핵심 요구, \code{multiset} 선택, one-occurrence erase, 복잡도, edge case와 T\_solve는 복원되었다. 다만 현재 회수 가능한 장기 기록에는 당시 제출 코드 전문과 sample 입력 전체가 verbatim으로 남아 있지 않아 아래에는 확인 가능한 실제 제출 구조를 수록한다. 확인 가능한 sample 출력은 \code{2 5}, \code{YES}, \code{2 5}, \code{NO}, \code{2 2}, \code{EMPTY} 순서였다.
\end{notebox}

\begin{answerbox}[학습자 제출 구조]
학습자는 \code{multiset<int>}를 사용했다.
\begin{lstlisting}
multiset<int> pool;

// ADD
pool.insert(x);

// CHECK
if (pool.find(x) != pool.end()) {
    cout << "YES\n";
}
else {
    cout << "NO\n";
}

// REMOVE one occurrence
auto it = pool.find(x);
if (it != pool.end()) {
    pool.erase(it);
}

// RANGE
if (pool.empty()) {
    cout << "EMPTY\n";
}
else {
    cout << *pool.begin() << ' '
         << *prev(pool.end()) << '\n';
}
\end{lstlisting}

\textbf{선택 이유}: 값만 저장하면 되고, 중복을 보존해야 하며, 동적 insert/remove와 정렬된 최소/최대가 필요하므로 \code{multiset}이 적절하다.

\textbf{시간복잡도}: 전체 $O(Q\log Q)$.

\textbf{공간복잡도}: $O(Q)$.

\textbf{주요 edge case}: 빈 상태의 \code{RANGE}, 존재하지 않는 값의 \code{REMOVE}, 중복 값, 음수 값.

\textbf{T\_solve}: 12분 32초.
\end{answerbox}

% =========================================================
% Q&A / remediation after Part F and Part G attempts
% =========================================================
\section{최종평가 이후 학습 질의응답과 보강}

\subsection{Q1. \code{set}과 \code{unordered\_set}의 가장 실용적인 차이는 무엇인가?}
둘 다 unique membership을 표현하지만, \code{set}은 ordered tree라서 lookup/insert/erase가 $O(\log N)$이고 sorted iteration과 boundary query가 가능하다. \code{unordered\_set}은 ordering을 제공하지 않지만 expected $O(1)$ lookup이 가능하다.

\subsection{Q2. Hash collision이 심하면 실제로 무엇이 달라지는가?}
많은 key가 같은 bucket에 몰리면 한 lookup에서 여러 candidate를 검사해야 한다. 따라서 \code{unordered\_map<int,...>}에서 저장된 distinct key 수를 $K$라 할 때 하나의 lookup이 worst-case $O(K)$까지 나빠질 수 있다.

\subsection{Q3. \code{vector<int> v = it->second;}와 \code{const vector<int>\& v = it->second;}의 차이는 무엇인가?}
첫 코드는 mapped vector의 길이가 $f$라면 $O(f)$ 시간과 $O(f)$ 추가 공간을 사용해 복사한다. 두 번째 코드는 기존 vector에 reference만 붙이므로 binding 자체는 $O(1)$이다.

\subsection{Q4. \code{find()} 결과를 왜 반드시 \code{end()}와 비교해야 하는가?}
찾지 못한 경우 iterator는 \code{end()}다. 이 iterator를 \code{it->second}처럼 dereference하면 잘못된 접근이다. 따라서 먼저 \code{it != mp.end()}를 확인한다.

\subsection{Q5. reverse index는 언제 필요한가?}
한 방향 mapping만 저장했는데 반대 방향 query가 반복되고, 매 query마다 전체 데이터를 훑게 된다면 reverse index를 유지하는 것이 유리할 수 있다. update 시 두 구조를 함께 갱신하는 비용과 공간을 지불하고 query를 빠르게 만든다.

\subsection{Q6. \code{map<int,set<int>>}를 순회할 때 왜 \code{const auto\&}가 안전한가?}
map의 element type은 \code{pair<const int,set<int>>}다. 직접 \code{pair<int,set<int>>}를 적으면 exact type이 달라 변환 임시 객체와 mapped container 복사가 생길 수 있다. \code{const auto\&}는 실제 element type을 그대로 reference한다.

\subsection{Q7. \code{multiset::find}, \code{erase(iterator)}, 최소/최대 접근 비용은?}
\begin{itemize}
\item \code{find(x)}: $O(\log N)$
\item \code{erase(iterator)}: amortized $O(1)$
\item \code{begin()}, \code{prev(end())}: $O(1)$ iterator 접근
\end{itemize}

여러 query에 걸쳐 성공적인 삭제 횟수의 총합은 이전 삽입 횟수를 넘을 수 없다는 global accounting도 사용할 수 있다.

\subsection{Q8. \code{lower\_bound}와 \code{upper\_bound}의 차이는?}
\code{lower\_bound(x)}는 $x$ 이상인 첫 값, \code{upper\_bound(x)}는 $x$보다 큰 첫 값이다. strictness가 다르므로 $x$ 자체가 container에 있을 때 결과가 달라진다.

\subsection{Q9. predecessor와 successor는 어떻게 얻는가?}
\begin{lstlisting}
auto right = s.lower_bound(x);
if (right != s.end()) {
    // *right: x 이상인 최소값
}
if (right != s.begin()) {
    auto left = prev(right);
    // *left: x보다 작은 최대값
}
\end{lstlisting}

\subsection{Q10. 왜 \code{end()}는 있는데 첫 원소 이전 iterator는 일반적으로 없는가?}
STL은 range를 $[\text{begin},\text{end})$로 모델링한다. \code{end()}는 순회를 종료하기 위한 경계 sentinel이다. forward range에서 \code{begin()} 이전 위치는 필요하지 않으며, 역방향 순회는 \code{rbegin()}, \code{rend()}라는 별도의 reverse range로 표현한다.

% =========================================================
% Part G
% =========================================================
\section{Part G --- Adaptive Extra Problem과 재평가 이력}

Part G에서는 최종적으로 통과한 두 문제만 남기는 대신, 실제 학습 과정에서 사용된 FAIL/VOID 문제와 remediation을 함께 기록한다. 이미 해설이나 핵심 관찰에 노출된 문제는 이후 독립 평가에 재사용하지 않는다.

\subsection{G-0. GPT-generated --- Symbol Balance Queries}
\begin{problembox}{초기 Extra Problem --- Symbol Balance Queries}
문자열 symbol별 정수 balance를 관리한다. 여러 update와 query를 처리하며, 각 symbol의 balance를 빠르게 갱신하고 질의해야 한다. 입력 문자열 전체의 총 길이를 $L$이라 둔다.

\textbf{당시 평가에서 확인된 핵심 조건}
\begin{itemize}
\item 문자열 key를 사용한다.
\item balance는 누적 갱신되며 \code{long long} 범위가 필요하다.
\item update 수를 $N$, query 수를 $M$, 입력되는 모든 문자열의 총 문자 수를 $L$로 둔다.
\item expected 시간복잡도와 공간복잡도를 설명한다.
\end{itemize}
\end{problembox}

\begin{notebox}[세션 복원 상태]
이 문제의 전체 statement, sample, 제출 코드 전문은 현재 회수 가능한 장기 기록에 verbatim으로 남아 있지 않다. 다만 실제 제출이 \code{unordered\_map<string,long long>} 기반이었고 코드 correctness가 맞았다는 점, 그리고 아래 complexity 설명 때문에 정식 FAIL이 된 사실은 평가 기록으로 확인된다.
\end{notebox}

\begin{answerbox}[당시 제출에서 확인되는 핵심]
\textbf{자료구조}: \code{unordered\_map<string,long long>}.

\textbf{당시 complexity 설명}: 문자열 hash 비용을 사실상 $O(1)$로 두어 전체를 $O(N+M)$으로 설명했다.

\textbf{당시 T\_solve}: 8분 11초.
\end{answerbox}

\begin{notebox}[당시 판정 --- FAIL / Complexity]
코드 자체와 \code{long long} 선택은 맞았지만, 문자열 key를 hash하고 입력하는 비용을 원소 개수만으로 계산했다. 총 문자열 길이 $L$을 포함하면 expected running time은
\[
O(L+N+M)
\]
처럼 써야 한다. 문자열 길이가 상수 상한으로 제한되면 $L=\Theta(N+M)$이므로 $O(N+M)$으로 단순화할 수 있지만, 당시 문제는 $L$을 명시적으로 모델링하도록 요구했다.
\end{notebox}

\subsection{G-0R. Remediation --- aggregate payload와 \code{operator[]}}
이 실패 뒤에는 $N$, $M$, update 문자열 총 길이 $L_u$, query 문자열 총 길이 $L_q$, 전체 $L=L_u+L_q$를 구분하는 연습을 했다.

핵심 결론은 다음과 같다.
\begin{itemize}
\item 문자열 hash/입력 비용까지 포함한 expected time은 $O(N+M+L)$이다.
\item 문자열 길이가 30 이하처럼 상수 제한이면 $L=\Theta(N+M)$이므로 $O(N+M)$도 asymptotically 맞다.
\item \code{mp[q]}는 없는 query key를 값 0으로 삽입할 수 있지만 \code{find(q)}는 상태를 바꾸지 않는다.
\item query를 별도 \code{vector}에 모아 둘 필요가 없다면 streaming 처리가 더 단순하다.
\end{itemize}

\subsection{G-1. External CP --- AtCoder ABC073 C Write and Erase}
\begin{problembox}{External CP --- Write and Erase}
처음에는 아무 숫자도 적혀 있지 않은 종이가 있다. $N$개의 정수 $A_1,\dots,A_N$을 순서대로 처리한다. 현재 종이에 $A_i$가 적혀 있으면 그 숫자를 지우고, 적혀 있지 않으면 적는다. 모든 처리가 끝난 뒤 종이에 남아 있는 서로 다른 숫자의 개수를 출력한다.

\textbf{제약 조건}
\begin{TextBlock}
1 <= N <= 100000
1 <= A_i <= 1000000000
\end{TextBlock}
\end{problembox}

\begin{notebox}[세션 복원 상태]
당시 제출 전문은 현재 장기 기록에 남아 있지 않지만, 실제 제출은 \code{unordered\_set<int>}에 대해 존재하면 erase, 없으면 insert하는 toggle 구조였고 \code{contains()}를 사용했다. T\_solve는 3분 55초였다.
\end{notebox}

\begin{notebox}[당시 판정 --- VOID]
직전 Symbol Balance Queries의 반복 complexity 오류로 formal-assessment pause가 이미 활성화된 뒤 제출되어 정식 PASS/FAIL evidence로 계산하지 않았다. 학습 모드 검토에서는 자료구조 방향과 알고리즘은 맞았다. \code{unordered\_set::contains()}는 C++20 기능이므로 정확히 C++17로 컴파일하면 사용할 수 없고, expected 전체 시간은 $O(N)$, hash collision까지 고려한 worst case는 $O(N^2)$이다.
\end{notebox}

\subsection{G-2. GPT-generated --- First Appearance Log}
\begin{problembox}{Fresh Reassessment --- First Appearance Log}
$N$개의 문자열이 순서대로 주어진다. 어떤 문자열이 처음 등장한 순간에만 그 위치와 문자열을 출력하고, 마지막에는 서로 다른 문자열의 총 개수를 출력한다. 입력 문자열들의 총 길이를 $L$이라 둔다.

\textbf{당시 요구사항}
\begin{itemize}
\item C++17 이상
\item 자료구조 선택 이유
\item expected 전체 시간복잡도와 worst-case 전체 시간복잡도
\item 공간복잡도
\item $N$과 $L$의 역할 설명
\item edge case와 T\_solve
\end{itemize}
\end{problembox}

\begin{answerbox}[당시 제출에서 보존된 코드 구조]
\begin{lstlisting}
unordered_set<string> apr;

// 각 입력 symbol에 대하여
if (apr.find(symbol) == apr.end()) {
    cout << pos << " " << symbol << '\n';
    apr.insert(symbol);
}
\end{lstlisting}

\textbf{expected 시간복잡도}: $O(L+N)$.

\textbf{공간복잡도}: 저장되는 문자열 payload를 포함해 $O(L)$.

\textbf{T\_solve}: 14분 2초.
\end{answerbox}

\begin{notebox}[당시 판정 --- FAIL / Complexity]
코드와 expected complexity는 맞았지만 요구된 worst-case complexity를 제출하지 않았다. 또한 비교 설명에서 unsorted \code{vector} membership을 $O(N\log N)$이라고 적었는데, 한 번의 선형 탐색은 $O(N)$이고 반복 query에서는 $O(N^2)$까지 갈 수 있다. \code{unordered\_set}도 collision이 심하면 operation 하나가 저장 key 수에 대해 선형으로 나빠질 수 있으므로 전체 worst case를 별도로 설명해야 했다.
\end{notebox}

\subsection{G-3. External CP --- AtCoder ABC235 C The Kth Time Query}
\begin{problembox}{Fresh Reassessment --- The Kth Time Query}
길이 $N$의 수열 $A=(a_1,\dots,a_N)$가 있다. 각 query $(x_i,k_i)$에 대해 수열을 앞에서부터 보았을 때 값 $x_i$가 $k_i$번째로 등장하는 위치를 출력한다. 그런 위치가 없으면 $-1$을 출력한다.

\textbf{제약 조건}
\begin{TextBlock}
1 <= N <= 200000
1 <= Q <= 200000
0 <= a_i, x_i <= 1000000000
1 <= k_i <= N
\end{TextBlock}

\textbf{예제 입력}
\begin{TextBlock}
6 8
1 1 2 3 1 2
1 1
1 2
1 3
1 4
2 1
2 2
2 3
4 1
\end{TextBlock}

\textbf{예제 출력}
\begin{TextBlock}
1
2
5
-1
3
6
-1
-1
\end{TextBlock}
\end{problembox}

\begin{answerbox}[당시 제출 코드 --- 세션 기록에서 복원]
\begin{lstlisting}
unordered_map<int, vector<int>> A;

for (int i = 0; i < N; ++i) {
    int Ai;
    cin >> Ai;
    A[Ai].push_back(i);
}

for (int i = 0; i < Q; ++i) {
    int xi, ki;
    cin >> xi >> ki;

    auto it = A.find(xi);
    vector<int> v = it->second;

    if (v.size() >= ki) {
        cout << v[ki - 1] << '\n';
    }
    else {
        cout << -1 << '\n';
    }
}
\end{lstlisting}

\textbf{자료구조}: \code{unordered\_map<int,vector<int>>}.

\textbf{당시 설명}: 전체를 expected $O(N+Q)$, 공간 $O(N)$으로 보았다.

\textbf{T\_solve}: 24분 9초.
\end{answerbox}

\begin{notebox}[당시 판정 --- FAIL / Correctness]
자료구조 방향은 맞았지만 세 가지 문제가 있었다.
\begin{enumerate}
\item 공식 출력은 1-based 위치인데 \code{i}를 그대로 저장해 0-based 위치를 출력했다.
\item \code{it == A.end()}일 수 있는데 확인 전에 \code{it->second}를 dereference했다.
\item \code{vector<int> v = it->second;}가 query마다 위치 vector 전체를 복사하므로 실제 complexity가 제출 설명과 달라질 수 있다.
\end{enumerate}
읽기만 한다면 \code{const vector<int>\&}로 reference를 붙이고, \code{end()}를 먼저 검사해야 한다.
\end{notebox}

\subsection{G-3R. Remediation --- iterator validity, indexing, copy/reference}
이후 다음 항목을 학습 모드에서 다시 확인했다.
\begin{itemize}
\item \code{find()} 결과가 \code{end()}인지 확인하기 전에는 dereference하지 않는다.
\item 출력 규칙이 1-based라면 저장 시 \code{i+1}을 넣거나 출력할 때 보정한다.
\item 길이 $f$의 \code{vector}를 값으로 복사하면 $O(f)$ 시간과 $O(f)$ 추가 공간이 든다.
\item \code{const vector<int>\& v = it->second;}의 reference binding 자체는 $O(1)$이다.
\item 정수 key \code{unordered\_map}의 lookup은 expected $O(1)$, distinct key 수가 $K$일 때 worst case $O(K)$이다.
\end{itemize}

\subsection{G-4. GPT-generated --- Active ID Registry}
\begin{problembox}{Fresh Reassessment --- Active ID Registry}
처음에는 활성화된 ID가 없다. $Q$개의 명령을 처리한다.
\begin{itemize}
\item \code{1 x}: ID \code{x}를 활성화한다. 이미 활성화되어 있으면 아무 일도 하지 않는다.
\item \code{2 x}: ID \code{x}를 비활성화한다. 활성화되어 있지 않으면 아무 일도 하지 않는다.
\item \code{3 x}: 현재 활성화되어 있으면 \code{1}, 아니면 \code{0}을 출력한다.
\item \code{4}: 현재 활성 ID의 개수를 출력한다.
\end{itemize}

\textbf{구현 및 제출 조건}
\begin{itemize}
\item $1\le Q\le200000$
\item $0\le x\le10^9$
\item C++17 이상
\item expected 전체 시간복잡도와 worst-case 전체 시간복잡도를 모두 설명한다.
\item 공간복잡도, numeric safety, edge case, T\_solve를 제출한다.
\end{itemize}
\end{problembox}

\begin{answerbox}[학습자 제출 답안]
\begin{lstlisting}
#include <iostream>
#include <string>
#include <algorithm>
#include <vector>
#include <map>
#include <set>
#include <unordered_map>
#include <unordered_set>

using namespace std;

int main()
{
    int Q = 0;
    cin >> Q;
    unordered_map<int, int> M = {};
    int act = 0;
    for (int i = 0; i < Q; i += 1) {
        int O1 = 0;
        int O2 = 0;
        int x = -1;
        cin >> O1;
        if (O1 == 1) {
            cin >> x;
            if (M[x] == 0) {
                act += 1;
                M[x] = 1;
            }
        }
        else if (O1 == 2) {
            cin >> x;
            if (M[x] == 1) {
                act += -1;
                M[x] = 0;
            }
        }
        else if (O1 == 3) {
            cin >> x;
            cout << M[x] << '\n';
        }
        else if (O1 == 4) {
            cout << act << '\n';
        }
    }
}
\end{lstlisting}

\textbf{자료구조}: \code{int}, \code{unordered\_map}.

\textbf{expected 시간복잡도}: $O(Q)$.

\textbf{worst-case 시간복잡도}: $O(Q^2)$.

\textbf{공간복잡도}: $O(Q)$.

\textbf{numeric safety}: 모든 정수 입력과 counter가 \code{int} 범위 안이다.

\textbf{edge case}: $Q=1$, \code{3 12}이면 \code{0}.

\textbf{T\_solve}: 17분 8초.
\end{answerbox}

\begin{notebox}[당시 판정 --- PASS]
Hash container의 expected/worst-case를 독립적으로 정확히 적용했다. 이 문제는 membership만 필요하므로 \code{unordered\_set<int>}가 더 직접적일 수 있지만, 제출한 \code{unordered\_map<int,int>}도 correctness를 만족한다. \code{M[x]}가 없는 key를 값 0으로 삽입할 수 있다는 점도 공간 분석에서 인식해야 한다.
\end{notebox}

\subsection{G-5. External CP --- AtCoder ABC298 C Cards Query Problem}
\begin{problembox}{External CP Reassessment --- Cards Query Problem}
$1$부터 $N$까지 번호가 붙은 빈 상자들이 있다. 다음 query를 처리한다.
\begin{itemize}
\item \code{1 i j}: 숫자 $i$가 적힌 카드 한 장을 상자 $j$에 넣는다.
\item \code{2 i}: 상자 $i$ 안의 카드 숫자를 오름차순으로 모두 출력한다. 중복도 모두 출력한다.
\item \code{3 i}: 숫자 $i$가 적힌 카드를 하나 이상 포함하는 상자 번호를 오름차순으로 출력한다. 같은 상자는 한 번만 출력한다.
\end{itemize}

\textbf{제약 조건}
\begin{TextBlock}
1 <= N, Q <= 200000
1 <= card number <= 200000
1 <= box number <= N
전체 출력 숫자 수 <= 200000
\end{TextBlock}

\textbf{예제 입력}
\begin{TextBlock}
5
8
1 1 1
1 2 4
1 1 4
2 4
1 1 4
2 4
3 1
3 2
\end{TextBlock}

\textbf{예제 출력}
\begin{TextBlock}
1 2
1 1 2
1 4
4
\end{TextBlock}
\end{problembox}

\begin{answerbox}[학습자 제출 답안]
\begin{lstlisting}
#include <iostream>
#include <string>
#include <algorithm>
#include <vector>
#include <map>
#include <set>
#include <unordered_map>
#include <unordered_set>

using namespace std;

int main()
{
    int N = 0;
    cin >> N;
    int Q = 0;
    cin >> Q;
    map<int, multiset<int>> boxes = {};
    for (int k = 0; k < Q; k += 1) {
        int O1 = 0;
        int O2 = 0;
        int O3 = 0;
        cin >> O1;
        if (O1 == 1) {
            cin >> O2 >> O3;
            boxes[O3].insert(O2);
        }
        else if (O1 == 2) {
            cin >> O2;
            const multiset<int>& O2thbox = boxes[O2];
            for (const int& i : O2thbox) {
                cout << i << " ";
            }
            cout << '\n';
        }
        else if (O1 == 3) {
            cin >> O2;
            for (const pair<int, multiset<int>>& p : boxes) {
                if (p.second.contains(O2)) {
                    cout << p.first << " ";
                }
            }
            cout << '\n';
        }
    }
}
\end{lstlisting}

\textbf{T\_solve}: 31분 27초.
\end{answerbox}

\begin{notebox}[당시 판정 --- FAIL / Complexity]
상자 내부에 중복 카드를 정렬 저장하기 위한 \code{multiset} 선택은 맞았다. 그러나 type 3에서 card $\to$ boxes reverse index가 없어 모든 상자를 순회했다. 또한 \code{map<int,multiset<int>>::value\_type}은 \code{pair<const int,multiset<int>>}이므로 \code{const pair<int,multiset<int>>\&}를 쓰면 변환 임시 객체가 생기며 큰 \code{multiset} 복사가 발생할 수 있다. 이 경우 \code{const auto\&}가 안전하다. type 2의 이미 정렬된 \code{multiset} 순회 출력은 $c$개 출력에 $O(c)$이다.
\end{notebox}

\subsection{G-5R. Remediation --- forward/reverse index와 \code{map::value\_type}}
양방향 query가 반복되면 다음 두 index를 함께 유지하는 사고를 연습했다.
\begin{lstlisting}
map<int, multiset<int>> box_to_cards;
map<int, set<int>> card_to_boxes;
\end{lstlisting}
전자는 중복 카드 보존, 후자는 같은 box를 한 번만 출력하기 위한 unique ordered set이다. checkpoint에서는 \code{user -> groups}만 있을 때 reverse query가 비효율적임을 설명하고 \code{group -> users}를 제안했으며, \code{map<int,set<int>>} 순회에는 \code{const auto\&}를 선택했다.

\subsection{G-6. External CP --- AtCoder ABC253 C Max - Min Query}
\begin{problembox}{External CP Reassessment --- Max - Min Query}
정수 multiset $S$가 처음에는 비어 있다.
\begin{itemize}
\item \code{1 x}: $x$ 하나를 추가한다.
\item \code{2 x c}: $x$를 최대 $c$개 삭제한다.
\item \code{3}: 현재 최댓값에서 최솟값을 뺀 값을 출력한다. 이 query에서 $S$는 비어 있지 않다.
\end{itemize}

\textbf{제약 조건}
\begin{TextBlock}
1 <= Q <= 200000
0 <= x <= 1000000000
1 <= c <= Q
\end{TextBlock}

\textbf{예제 입력}
\begin{TextBlock}
8
1 3
1 2
3
1 2
1 7
3
2 2 3
3
\end{TextBlock}

\textbf{예제 출력}
\begin{TextBlock}
1
5
4
\end{TextBlock}
\end{problembox}

\begin{answerbox}[학습자 제출 답안]
\begin{lstlisting}
#include <iostream>
#include <string>
#include <algorithm>
#include <vector>
#include <map>
#include <set>
#include <unordered_map>
#include <unordered_set>

using namespace std;

int main()
{
    int Q = 0;
    cin >> Q;
    multiset<int> S = {};
    for (int i = 0; i < Q; i += 1) {
        int order1 = 0;
        int x = 0;
        int c = 0;
        cin >> order1;
        if (order1 == 1) {
            cin >> x;
            S.insert(x);
        }
        else if (order1 == 2) {
            cin >> x >> c;
            for (int j = 0; j < c; j += 1) {
                auto it = S.find(x);
                if (it == S.end()) {
                    break;
                }
                else {
                    S.erase(it);
                }
            }
        }
        else if (order1 == 3) {
            cout << *prev(S.end()) - *S.begin() << '\n';
        }
    }
}
\end{lstlisting}

\textbf{T\_solve}: 8분 50초.
\end{answerbox}

\begin{notebox}[당시 판정 --- FAIL / Complexity]
코드와 \code{multiset} 선택은 맞았지만 complexity 설명이 부정확했다. \code{S.begin()}과 nonempty \code{prev(S.end())} 접근은 $O(1)$이고, \code{find(x)}는 $O(\log |S|)$, \code{erase(iterator)}는 amortized $O(1)$이다. 모든 type 2 query에서 성공한 삭제 횟수의 총합은 이전 insert 횟수를 넘을 수 없으므로 전체 프로그램은 $O(Q\log Q)$로 정리할 수 있다.
\end{notebox}

\subsection{G-6R. Remediation --- operation cost와 global accounting}
checkpoint에서 다음을 다시 확인했다.
\begin{itemize}
\item \code{find(x)}: $O(\log |S|)$
\item \code{erase(iterator)}: amortized $O(1)$
\item \code{begin()}, \code{prev(end())}: $O(1)$
\item 전체 성공 삭제 횟수 $\le$ 전체 삽입 횟수
\item 따라서 총 query 비용은 $O(Q\log Q)$
\end{itemize}

\subsection{G-7. External CP --- AtCoder ABC217 D Cutting Woods}
\begin{problembox}{External CP Reassessment --- Cutting Woods}
길이 $L$인 목재가 있고 처음에는 잘려 있지 않다.
\begin{itemize}
\item \code{1 x}: 위치 $x$에서 목재를 자른다.
\item \code{2 x}: 현재 위치 $x$를 포함하는 목재 조각의 길이를 출력한다.
\end{itemize}
query 시점에는 $x$가 기존 절단점이 아님이 보장된다.

\textbf{제약 조건}
\begin{TextBlock}
1 <= L <= 1000000000
1 <= Q <= 200000
1 <= x <= L-1
\end{TextBlock}

\textbf{예제 입력}
\begin{TextBlock}
5 3
2 2
1 3
2 2
\end{TextBlock}

\textbf{예제 출력}
\begin{TextBlock}
5
3
\end{TextBlock}
\end{problembox}

\begin{answerbox}[학습자 제출 답안]
\begin{lstlisting}
#include <iostream>
#include <string>
#include <algorithm>
#include <vector>
#include <map>
#include <set>
#include <unordered_map>
#include <unordered_set>

using namespace std;

int main()
{
    int L = 0;
    cin >> L;
    int Q;
    cin >> Q;
    set<int> wood = {0, L};
    int spoint = 0;
    for (int i = 0; i < Q; i += 1) {
        int q = 0;
        int x = 0;
        cin >> q >> x;
        if (q == 1) {
            wood.insert(x);
        }
        else {
            for (const auto& p : wood) {
                if (p != 0) {
                    if (spoint < x && x < p) {
                        cout << p - spoint << '\n';
                    }
                    spoint = p;
                }
            }
            spoint = 0;
        }
    }
}
\end{lstlisting}

\textbf{T\_solve}: 22분 31초.
\end{answerbox}

\begin{notebox}[당시 판정 --- FAIL / Complexity]
\code{set<int>} 선택과 코드 correctness는 맞았지만 type 2에서 set 전체를 선형 순회해 $O(|wood|)$을 사용했다. ordered set의 핵심 장점인 \code{lower\_bound(x)}를 이용하면 오른쪽 절단점을 $O(\log |wood|)$에 찾고 \code{prev()}로 왼쪽 절단점을 얻을 수 있으므로 전체를 $O(Q\log Q)$로 만들 수 있다.
\end{notebox}

\subsection{G-7R. Remediation --- \code{lower\_bound}, \code{upper\_bound}, predecessor/successor}
\code{set<int> s = \{3,7,12,20\}}에서 다음을 확인했다.
\begin{itemize}
\item \code{lower\_bound(7)} $\to 7$
\item \code{upper\_bound(7)} $\to 12$
\item \code{lower\_bound(9)} $\to 12$
\item \code{upper\_bound(20)} $\to$ \code{end()}
\end{itemize}
$x=10$일 때 왼쪽 값은 \code{prev(s.lower\_bound(x))}, 오른쪽 값은 \code{s.lower\_bound(x)}이며, \code{prev(s.lower\_bound(x))}가 안전하려면 \code{s.lower\_bound(x) != s.begin()}이어야 한다.

\subsection{G-8. External CP --- AtCoder ABC241 D Sequence Query}
\begin{problembox}{Fresh External Reassessment --- Sequence Query}
처음에는 빈 multiset $A$가 있다.
\begin{itemize}
\item \code{1 x}: $x$를 하나 추가한다.
\item \code{2 x k}: $x$ 이하 원소 중 큰 쪽에서 $k$번째 값을 출력하고, 없으면 $-1$.
\item \code{3 x k}: $x$ 이상 원소 중 작은 쪽에서 $k$번째 값을 출력하고, 없으면 $-1$.
\end{itemize}

\textbf{제약 조건}
\begin{TextBlock}
1 <= Q <= 200000
1 <= x <= 1000000000000000000
1 <= k <= 5
\end{TextBlock}

\textbf{예제 입력}
\begin{TextBlock}
11
1 20
1 10
1 30
1 20
3 15 1
3 15 2
3 15 3
3 15 4
2 100 5
1 1
2 100 5
\end{TextBlock}

\textbf{예제 출력}
\begin{TextBlock}
20
20
30
-1
-1
1
\end{TextBlock}
\end{problembox}

\begin{answerbox}[학습자 제출 답안]
\begin{lstlisting}
#include <iostream>
#include <string>
#include <algorithm>
#include <vector>
#include <map>
#include <set>
#include <unordered_map>
#include <unordered_set>

using namespace std;

int main()
{
    int Q = 0;
    cin >> Q;
    multiset<long long> A = {};
    for (int i = 0; i < Q; i += 1) {
        int q = 0;
        long long x = 0;
        int k = 0;
        cin >> q;
        if (q == 1) {
            cin >> x;
            A.insert(x);
        }
        else if (q == 2) {
            cin >> x >> k;
            auto it = A.upper_bound(x);
            auto it2 = A.begin();
            if (it == A.begin()) {
                cout << -1 << '\n';
            }
            else {
                for (int j = 1; j <= k; j += 1) {
                    it--;
                    if (it == it2 && j < k) {
                        cout << -1 << '\n';
                        break;
                    }
                    else if (j == k) {
                        cout << *it << '\n';
                    }
                }
            }
        }
        else if (q == 3) {
            cin >> x >> k;
            auto it = A.lower_bound(x);
            for (int j = 1; j <= k; j += 1) {
                if (j == 1) {
                }
                else {
                    if (it == A.end()) {
                        break;
                    }
                    else {
                        it = next(it);
                    }
                }
            }
            if (it != A.end()) {
                cout << *it << '\n';
            }
            else {
                cout << -1 << '\n';
            }
        }
    }
}
\end{lstlisting}

\textbf{자료구조}: \code{multiset<long long>}.

\textbf{q1}: $O(\log |A|)$.

\textbf{q2, q3}: bound search가 $O(\log |A|)$이고 $k\le5$이므로 iterator 이동은 상수 횟수다. 따라서 각각 $O(\log |A|)$.

\textbf{전체 시간복잡도}: $O(Q\log Q)$.

\textbf{공간복잡도}: $O(Q)$.

\textbf{numeric safety}: $x\le10^{18}$이므로 \code{long long}; $Q$, $k$, loop counter는 \code{int}로 충분하다.

\textbf{edge case}: 후보 원소가 $k$개보다 적어 $-1$을 출력하는 경우.

\textbf{T\_solve}: 19분 12초.
\end{answerbox}

\begin{notebox}[당시 판정 --- PASS]
이전 remediation의 핵심이 실제 외부 문제로 전이되었다. 중복 보존을 위해 \code{multiset<long long>}을 선택했고, type 2에서 \code{upper\_bound}, type 3에서 \code{lower\_bound}를 사용했으며 iterator boundary와 $k\le5$를 정확히 처리했다. 공식 sample과 randomized differential test를 통과했다.
\end{notebox}

% =========================================================
% Part H
% =========================================================
\section{Part H --- 숙달 점검과 복습}

\subsection{H1. 실제 평가 흐름 요약}
\begin{longtable}{p{4.3cm}p{2.0cm}p{2.0cm}p{6.0cm}}
\toprule
\textbf{문제} & \textbf{결과} & \textbf{T\_solve} & \textbf{핵심 학습점} \\
\midrule
Ticker Directory & PASS & 19:51 & ordered key $\to$ value, dynamic update, lexicographic FIRST \\
Live Value Pool & PASS & 12:32 & duplicate-preserving \code{multiset}, one-occurrence erase \\
Symbol Balance Queries & FAIL & 8:11 & 문자열 payload $L$ 누락, complexity accounting \\
ABC073 C Write and Erase & VOID & 3:55 & pause 중 제출; unordered membership toggle, C++20 \code{contains} \\
First Appearance Log & FAIL & 14:02 & worst-case complexity 누락 \\
ABC235 C The Kth Time Query & FAIL & 24:09 & 1-based indexing, \code{end()} validity, vector copy \\
Active ID Registry & PASS & 17:08 & unordered expected/worst-case complexity \\
ABC298 C Cards Query Problem & FAIL & 31:27 & reverse index, \code{map::value\_type}, output cost \\
ABC253 C Max - Min Query & FAIL & 8:50 & ordered-container operation cost, global accounting \\
ABC217 D Cutting Woods & FAIL & 22:31 & set을 골라도 핵심 \code{lower\_bound}를 써야 함 \\
ABC241 D Sequence Query & PASS & 19:12 & lower/upper bound, predecessor/successor, boundary safety \\
\bottomrule
\end{longtable}

\subsection{H2. 개념 자기점검}
다음 질문에 자료 없이 답할 수 있는지 확인하라.
\begin{enumerate}
\item membership만 필요할 때 map보다 set 계열이 자연스러운 이유는 무엇인가?
\item \code{map}과 \code{unordered\_map}의 ordering 및 expected/worst-case complexity 차이는 무엇인가?
\item \code{set}과 \code{multiset}의 중복 처리 차이는 무엇인가?
\item \code{operator[]}가 lookup만 하는 연산이 아닌 이유는 무엇인가?
\item \code{find()}가 실패했을 때 반환하는 iterator는 무엇인가?
\item \code{erase(iterator)}와 \code{erase(key)}를 구분해야 하는 이유는 무엇인가?
\item 문자열 key 문제에서 총 문자 수 $L$을 따로 정의하는 이유는 무엇인가?
\item 어떤 query 패턴에서 reverse index가 필요한가?
\item \code{map<K,V>::value\_type}의 key가 \code{const K}인 이유는 무엇인가?
\item \code{lower\_bound(x)}와 \code{upper\_bound(x)}의 strictness 차이는 무엇인가?
\item $x$보다 작은 최대값을 ordered set에서 어떻게 찾는가?
\item 왜 \code{prev(begin())}은 invalid이고 nonempty container의 \code{prev(end())}은 valid한가?
\end{enumerate}

\subsection{H3. 구현 자기점검}
\begin{checklistbox}[S1-A 코드 작성 후 체크리스트]
\begin{itemize}
\item[$\square$] key만 필요한지 key $\to$ value가 필요한지 구분했는가?
\item[$\square$] ordering이 실제로 필요한가?
\item[$\square$] 중복을 보존해야 하는가?
\item[$\square$] \code{operator[]}가 의도치 않은 key를 삽입하지 않는가?
\item[$\square$] \code{find()} 결과를 \code{end()} 확인 전에 dereference하지 않았는가?
\item[$\square$] 0-based / 1-based 요구를 정확히 맞췄는가?
\item[$\square$] 큰 mapped container를 값으로 복사하지 않는가?
\item[$\square$] reverse query를 매번 전체 scan하고 있지 않은가?
\item[$\square$] ordered container를 골랐다면 \code{lower\_bound}/\code{upper\_bound} 같은 핵심 연산을 활용하고 있는가?
\item[$\square$] \code{prev(it)} 전에 \code{it != begin()}을 확인했는가?
\item[$\square$] \code{*it} 전에 \code{it != end()}을 확인했는가?
\item[$\square$] unordered container의 complexity를 expected와 worst-case로 구분할 필요가 있는가?
\item[$\square$] 문자열 key의 문자 처리 비용을 무시해도 되는 제한인지 확인했는가?
\item[$\square$] 출력 개수 자체의 $O(K)$ 비용을 포함했는가?
\item[$\square$] 여러 query를 합산할 때 total insertion/deletion 같은 global bound를 사용할 수 있는가?
\end{itemize}
\end{checklistbox}

\subsection{H4. S1-A에서 형성해야 할 핵심 사고}
문제를 읽자마자 특정 container 이름을 떠올리는 것보다 다음 순서로 판단한다.
\begin{enumerate}
\item 무엇을 key로 찾는가?
\item value 상태도 함께 저장해야 하는가?
\item 순서가 필요한가?
\item 중복을 보존해야 하는가?
\item insert / erase가 동적으로 반복되는가?
\item 최소/최대 또는 경계 이웃을 찾아야 하는가?
\item query 방향이 하나인가, 양방향인가?
\item 요구 시간복잡도에서 선형 scan이 허용되는가?
\item 선택한 container의 \emph{핵심 연산}을 실제로 사용하고 있는가?
\end{enumerate}

\subsection{H5. 현재 복습 상태와 다음 Unit}
즉시 재평가에서는 ordered/unordered 선택과 boundary query가 최종적으로 독립 전이되었지만, 여러 External CP 시도에서 complexity 설명과 operation selection 오류가 반복되었다. 따라서 즉시 prerequisite block은 해제하되 delayed/mixed review에서 한 번 더 확인할 가치가 있다.

다음 Learning Unit인 S1-B에서는 Sorting \& Bounds를 더 일반적으로 다룬다.
\begin{itemize}
\item 정렬된 \code{vector}에서의 \code{lower\_bound}, \code{upper\_bound}
\item binary search의 경계 해석
\item custom comparator와 정렬 기준
\item ordered container와 정렬된 sequence의 trade-off
\end{itemize}


\subsection{H6. 최종 Mastery Record}

\begin{itemize}
\item \textbf{Learning Unit}: S1-A --- Associative Containers
\item \textbf{Part E}: PASS --- Ticker Directory
\item \textbf{Part F}: PASS --- Live Value Pool
\item \textbf{Part G}: 완료 --- GPT-generated fresh reassessment \textit{Active ID Registry} PASS, 최종 External CP reassessment AtCoder ABC241 D \textit{Sequence Query} PASS
\item \textbf{Capability}: L4 --- Selection
\item \textbf{Confidence}: Provisional
\item \textbf{Evidence Context}: Immediate
\item \textbf{Unit Coverage}: Sufficient for Provisional; immediate progression gate satisfied
\item \textbf{Review Debt}: Open --- Medium
\item \textbf{Immediate Retest}: No
\item \textbf{Next Review}: delayed/mixed assessment에서 ordered associative-container 선택, complexity accounting, reverse-index 선택 경계를 재검증
\end{itemize}

\paragraph{Core Decision Boundary Coverage.}
\begin{itemize}
\item key $\to$ value + ordering 필요: \textbf{Immediate PASS} --- Ticker Directory에서 \code{map} 선택.
\item duplicate-preserving ordered dynamic state: \textbf{Immediate PASS} --- Live Value Pool과 ABC241 D에서 \code{multiset} 선택.
\item ordering 불필요 + hash 기반 membership/state 및 expected/worst-case complexity: \textbf{Immediate PASS} --- Active ID Registry.
\item predecessor/successor와 ordered boundary query: \textbf{Immediate PASS} --- ABC217 D 실패 후 remediation을 거쳐 ABC241 D에서 독립 PASS.
\item forward/reverse index 선택: \textbf{Not Yet Independently Confirmed} --- ABC298 C 실패 후 remediation은 통과했지만 fresh independent formal PASS가 아직 없다.
\item static sorted sequence vs dynamic ordered container: \textbf{Not Yet Independently Confirmed} --- 다음 Unit S1-B와 delayed/mixed evidence에서 보완한다.
\end{itemize}

\paragraph{종료 판정.}
S1-A의 Immediate Coverage Floor와 progression gate는 충족되었다. 따라서 다음 Learning Unit인 S1-B --- Sorting \& Bounds로 진행한다. 다만 여러 External CP 시도에서 operation-selection 및 complexity 설명 오류가 반복되었으므로 Confidence는 Provisional로 유지하고 Review Debt Medium을 닫지 않는다. Confirmed 승격은 주제 cue가 제거된 delayed/mixed assessment에서 Core Decision Boundary가 다시 독립적으로 검증된 뒤 판단한다.
